\clearpage
\begin{trapbox}
\textbf{\textcolor{CTrap}{易错一：只知道“不下降”，却当成“严格递增”}}
\begin{center}\begin{tikzpicture}[font=\small,>=stealth,line join=round]
\draw[gray!40,rounded corners=3pt] (0,0) rectangle (14.2,8.3);\node at (7.1,7.700000000000001) {不下降可以一直平：无法保留严格不等号};
\begin{scope}[shift={(4.1,3.9499999999999997)},x=1.0cm,y=1.15cm]\draw[->,gray] (-3,0)--(3,0) node[right]{$x$};\draw[->,gray] (0,-2)--(0,2) node[above]{$y$};\draw[blue,thick,solid,domain=-2.5:2.5,samples=180] plot(\x,{0});\fill[orange!85!black] (1,0) circle(2.2pt);\end{scope}\node[anchor=west,align=left,text width=5.0cm] at (8.4,4.0) {$f(x)=0$\\[9pt]奇且偶，处处不下降\\[9pt]$a=b=1$ 时 $a+b=2>0$\\[9pt]但 $f(a)+f(b)=0$};
\node[text width=12.6cm,align=center] at (7.1,.65) {常值函数给出反例；严格单调才能把函数值比较与输入比较双向转换。};
\end{tikzpicture}\end{center}



\textbf{错误做法：}函数在右半边“不下降”，就写
\[
 f(a)+f(b)>0\iff a+b>0.
\]
\textbf{为什么错？}取 $f(x)=0$。它既是奇函数又是偶函数，而且从左到右都不会下降。取 $a=b=1$，有 $a+b=2>0$，但 $f(1)+f(1)=0$，并不大于 $0$。

\textbf{正确做法：}这里必须有\textbf{严格}递增或\textbf{严格}递减，才能把函数值的严格大小关系和输入的严格大小关系互相转换。

\medskip

\end{trapbox}
\clearpage
\begin{trapbox}
\textbf{\textcolor{CTrap}{易错二：以为奇函数的右半边递增，就自动有 $f(x)>0$（忘了 $0$）}}
\begin{center}\begin{tikzpicture}[font=\small,>=stealth,line join=round]
\draw[gray!40,rounded corners=3pt] (0,0) rectangle (14.2,8.3);\node at (7.1,7.700000000000001) {缺少零点的两段定义域，不能拼成整体递增};
\begin{scope}[shift={(4.1,3.9499999999999997)},x=0.8108108108108107cm,y=0.019166666666666665cm]\draw[->,gray] (-3.7,0)--(3.7,0) node[right]{$x$};\draw[->,gray] (0,-120)--(0,120) node[above]{$y$};\draw[blue,thick,solid,domain=-3:-1,samples=180] plot(\x,{\x+100});\draw[blue,thick,solid,domain=1:3,samples=180] plot(\x,{\x-100});\fill[orange!85!black] (-3,97) circle(2.2pt);\fill[orange!85!black] (-1,99) circle(2.2pt);\fill[orange!85!black] (1,-99) circle(2.2pt);\fill[orange!85!black] (3,-97) circle(2.2pt);\end{scope}\node[anchor=west,align=left,text width=5.0cm] at (8.4,4.0) {$D=[-3,-1]\cup[1,3]$\\[9pt]左支 $x+100$\\[9pt]右支 $x-100$\\[9pt]$f(2)=-98<0$\\[9pt]虽奇且右支严格递增};
\node[text width=12.6cm,align=center] at (7.1,.65) {两段之间没有曲线；不能跨过定义域缺口，把局部递增当成整体递增。};
\end{tikzpicture}\end{center}



\textbf{关键原因：}我们说右半边严格递增的奇函数在整个区间递增，用到了 $0$ 在定义域内，并且奇函数必有 $f(0)=0$。

\textbf{反例提醒：}若定义域不是以 $0$ 为中间位置的区间，而是 $[-3,-1]\cup[1,3]$，只知道右半边严格递增还不够。比如
\[
 f(x)=\begin{cases}x-100,&1\leqslant x\leqslant3,\\x+100,&-3\leqslant x\leqslant-1.\end{cases}
\]
它是奇函数，右半边严格递增；但取 $a=b=2$，有 $a+b=4>0$，却有 $f(2)+f(2)=-196<0$。

\textbf{正确做法：}本结论使用的定义域限定为 $\mathbb R$ 或一个包含 $0$、左右对称的区间；若题目给了不连在一起的定义域，应重新判断，不能照套。

\medskip

\end{trapbox}
\clearpage
\begin{trapbox}
\textbf{\textcolor{CTrap}{易错三：把奇函数规则套在偶函数上}}
\begin{center}\begin{tikzpicture}[font=\small,>=stealth,line join=round]
\draw[gray!40,rounded corners=3pt] (0,0) rectangle (14.2,8.3);\node at (7.1,7.700000000000001) {偶函数可以整体向下平移：和式符号随之改变};
\begin{scope}[shift={(4.1,5.866666666666667)},x=0.8571428571428571cm,y=0.03833333333333333cm]\draw[->,gray] (-3.5,0)--(3.5,0) node[right]{$x$};\draw[->,gray] (0,-110)--(0,10) node[above]{$y$};\draw[blue,thick,solid,domain=-3:3,samples=180] plot(\x,{(\x)^2-100});\fill[orange!85!black] (1,-99) circle(2.2pt);\fill[orange!85!black] (2,-96) circle(2.2pt);\end{scope}\node[anchor=west,align=left,text width=5.0cm] at (8.4,4.0) {$f(x)=x^2-100$\\[9pt]右半边严格递增\\[9pt]$a=1,\ b=2$，和为 $3$\\[9pt]$f(1)+f(2)=-195<0$};
\node[text width=12.6cm,align=center] at (7.1,.65) {左右同高控制两值的比较，不能仅由输入之和判断函数值之和的符号。};
\end{tikzpicture}\end{center}



\textbf{错误做法：}把 $f(a)+f(b)>0$ 直接变成 $a+b>0$，不管函数是奇是偶。

\textbf{反例：}取偶函数 $f(x)=x^2-100$，右半边严格递增。令 $a=1,b=2$，虽然 $a+b=3>0$，但
\[
 f(1)+f(2)=-99-96=-195<0.
\]
\textbf{正确做法：}奇函数的\textbf{和式}可变成比较 $a$ 与 $-b$；偶函数的\textbf{两个函数值比大小}，要比 $|a|$ 与 $|b|$。偶函数的和式正负，通常不能只由 $a+b$ 判断。

\medskip

\end{trapbox}
\clearpage
\begin{trapbox}
\textbf{\textcolor{CTrap}{易错四：偶函数直接比 $a$、$b$，忘了绝对值}}
\begin{center}\begin{tikzpicture}[font=\small,>=stealth,line join=round]
\draw[gray!40,rounded corners=3pt] (0,0) rectangle (14.2,8.3);\node at (7.1,7.700000000000001) {偶函数：先把负输入折到右边，再比距离};
\begin{scope}[shift={(4.1,2.003846153846154)},x=0.8108108108108107cm,y=0.3538461538461538cm]\draw[->,gray] (-3.7,0)--(3.7,0) node[right]{$x$};\draw[->,gray] (0,-1)--(0,12) node[above]{$y$};\draw[blue,thick,solid,domain=-3.2:3.2,samples=180] plot(\x,{(\x)^2});\fill[orange!85!black] (-3,9) circle(2.2pt);\fill[orange!85!black] (3,9) circle(2.2pt);\fill[orange!85!black] (2,4) circle(2.2pt);\draw[violet,dashed] (-3,9)--(3,9);\end{scope}\node[anchor=west,align=left,text width=5.0cm] at (8.4,4.0) {图取 $f(x)=x^2$\\[9pt]$a=-3,\ b=2$\\[9pt]$f(-3)=f(3)=9>4=f(2)$\\[9pt]比较 $|-3|=3>|2|=2$};
\node[text width=12.6cm,align=center] at (7.1,.65) {横坐标较小不等于函数值较小；真正比较的是到纵轴的水平距离。};
\end{tikzpicture}\end{center}



\textbf{反例：}取 $f(x)=x^2$，有 $-3<2$，却有 $f(-3)=9>4=f(2)$。

\textbf{正确做法：}把 $-3$ 和 $2$ 改看成距离 $3$ 和 $2$；偶函数右半边严格递增，则 $3>2$ 对应 $f(-3)>f(2)$。

\medskip

\end{trapbox}
\clearpage
\begin{trapbox}
\textbf{\textcolor{CTrap}{易错五：严格递减时，不等号没有反过来}}
\begin{center}\begin{tikzpicture}[font=\small,>=stealth,line join=round]
\draw[gray!40,rounded corners=3pt] (0,0) rectangle (14.2,8.3);\node at (7.1,7.700000000000001) {严格递减时：输入更小，函数值反而更大};
\begin{scope}[shift={(4.1,3.9499999999999997)},x=1.0cm,y=0.7666666666666666cm]\draw[->,gray] (-3,0)--(3,0) node[right]{$x$};\draw[->,gray] (0,-3)--(0,3) node[above]{$y$};\draw[blue,thick,solid,domain=-2.5:2.5,samples=180] plot(\x,{-\x});\fill[orange!85!black] (-2,2) circle(2.2pt);\fill[orange!85!black] (1,-1) circle(2.2pt);\end{scope}\node[anchor=west,align=left,text width=5.0cm] at (8.4,4.0) {奇函数 $f(x)=-x$\\[9pt]$a=-2,\ b=1$\\[9pt]$a+b=-1<0$\\[9pt]$f(-2)+f(1)=2-1=1>0$};
\node[text width=12.6cm,align=center] at (7.1,.65) {将输入比较的方向反过来；等号位置不受递增或递减影响。};
\end{tikzpicture}\end{center}



\textbf{看一个最容易核对的函数：}奇函数 $f(x)=-x$ 是严格递减的。取 $a=-2,b=1$，则 $a+b=-1<0$，却有
\[
 f(-2)+f(1)=2-1=1>0.
\]
\textbf{记住：}奇函数严格递减时，和为正对应 $a+b<0$；偶函数严格递减时，函数值较大对应距离 $|x|$ 较小。

\medskip

\end{trapbox}
\clearpage
\begin{trapbox}
\textbf{\textcolor{CTrap}{易错六：把两个数的结论强行推广到三个数}}
\begin{center}\begin{tikzpicture}[font=\small,>=stealth,line join=round]
\draw[gray!40,rounded corners=3pt] (0,0) rectangle (14.2,8.3);\node at (7.1,7.700000000000001) {三个输入的和为零，立方后的和却不为零};
\begin{scope}[shift={(3.6714285714285717,3.183333333333333)},x=1.4285714285714286cm,y=0.30666666666666664cm]\draw[->,gray] (-1.8,0)--(2.4,0) node[right]{$x$};\draw[->,gray] (0,-5)--(0,10) node[above]{$y$};\draw[blue,thick,solid,domain=-1.4:2.1,samples=180] plot(\x,{(\x)^3});\fill[orange!85!black] (-1,-1) circle(2.2pt);\fill[orange!85!black] (2,8) circle(2.2pt);\end{scope}\node[anchor=west,align=left,text width=5.0cm] at (8.4,4.0) {$f(x)=x^3$\\[9pt]输入 $2,-1,-1$\\[9pt]输入和 $2-1-1=0$\\[9pt]函数值和 $8-1-1=6$\\[9pt]$-1$ 这个输入使用两次};
\node[text width=12.6cm,align=center] at (7.1,.65) {曲线并非把加法原样搬到纵轴；两项结论不能直接推广到三项。};
\end{tikzpicture}\end{center}



\textbf{错误猜想：}“既然奇函数 $f(a)+f(b)$ 的正负由 $a+b$ 决定，三个数相加也一样吧？”

\textbf{反例：}取 $f(x)=x^3$，$a=2,b=-1,c=-1$，则
\[
 a+b+c=0,\qquad f(a)+f(b)+f(c)=8-1-1=6>0.
\]
\textbf{正确理解：}两个数之所以好处理，是奇函数能把其中一个 $f(b)$ 直接换成 $-f(-b)$，马上只剩下\textbf{两个函数值比较}。三个函数值相加一般做不到这一步。

\medskip

\end{trapbox}
\clearpage
\begin{trapbox}
\textbf{\textcolor{CTrap}{易错七：绝对值不等式只解了一半}}
\begin{center}\begin{tikzpicture}[font=\small,>=stealth,line join=round]
\draw[gray!40,rounded corners=3pt] (0,0) rectangle (14.2,8.8);\node at (7.1,8.200000000000001) {例2：绝对值大于1，两边的解都要保留};
\node[anchor=west] at (.6,6.8) {令 $u=2a-1$，解 $|u|>1$：};\draw[->,gray] (1,6.3)--(13,6.3);\draw[blue,very thick] (1.3,6.3)--(5,6.3);\draw[blue,very thick] (8,6.3)--(12,6.3);\draw[blue,fill=white] (5,6.3) circle(2.7pt);\node[below=7pt] at (5,6.3) {$-1$};\draw[blue,fill=white] (8,6.3) circle(2.7pt);\node[below=7pt] at (8,6.3) {$1$};\node[anchor=west] at (.6,5.0) {换回 $a$，解为 $a<0$ 或 $a>1$：};\draw[->,gray] (1,4.5)--(13,4.5);\draw[violet,very thick] (1.3,4.5)--(5,4.5);\draw[violet,very thick] (8,4.5)--(12,4.5);\draw[violet,fill=white] (5,4.5) circle(2.7pt);\node[below=7pt] at (5,4.5) {$0$};\draw[violet,fill=white] (8,4.5) circle(2.7pt);\node[below=7pt] at (8,4.5) {$1$};
\node[text width=12.6cm,align=center] at (7.1,.65) {各行是各自的数轴；空心端点表示不取等号，两条射线缺一不可。};
\end{tikzpicture}\end{center}



例如 $|2a-1|>1$ 必须分成 $2a-1>1$ \textbf{或} $2a-1<-1$，因此 $a>1$ \textbf{或} $a<0$。漏掉负数那一边，就丢掉了偶函数图像左侧对应的解。

\medskip

\end{trapbox}
\clearpage
\begin{trapbox}
\textbf{\textcolor{CTrap}{易错八：忘记函数的定义域}}
\begin{center}\begin{tikzpicture}[font=\small,>=stealth,line join=round]
\draw[gray!40,rounded corners=3pt] (0,0) rectangle (14.2,8.8);\node at (7.1,8.200000000000001) {例3：先求代数解，再保留可以代入的输入};
\node[anchor=west] at (.6,6.8) {比较得到 $t<-\frac12$：};\draw[->,gray] (1,6.3)--(13,6.3);\draw[blue,very thick] (1.3,6.3)--(6.5,6.3);\draw[blue,fill=white] (6.5,6.3) circle(2.7pt);\node[below=7pt] at (6.5,6.3) {$-\frac12$};\node[anchor=west] at (.6,5.0) {输入限制 $-3\le t\le2$：};\draw[->,gray] (1,4.5)--(13,4.5);\draw[violet,very thick] (3,4.5)--(10,4.5);\draw[violet,fill=violet] (3,4.5) circle(2.7pt);\node[below=7pt] at (3,4.5) {$-3$};\draw[violet,fill=violet] (10,4.5) circle(2.7pt);\node[below=7pt] at (10,4.5) {$2$};\node[anchor=west] at (.6,3.1999999999999997) {共同部分 $[-3,-\frac12)$：};\draw[->,gray] (1,2.6999999999999997)--(13,2.6999999999999997);\draw[orange!85!black,very thick] (3,2.6999999999999997)--(6.5,2.6999999999999997);\draw[orange!85!black,fill=orange!85!black] (3,2.6999999999999997) circle(2.7pt);\node[below=7pt] at (3,2.6999999999999997) {$-3$};\draw[orange!85!black,fill=white] (6.5,2.6999999999999997) circle(2.7pt);\node[below=7pt] at (6.5,2.6999999999999997) {$-\frac12$};
\node[text width=12.6cm,align=center] at (7.1,.65) {三行共用同一 $t$ 尺度；最左端可取，右端使和等于0而排除。};
\end{tikzpicture}\end{center}



若 $f$ 只定义在 $[-3,3]$，遇到 $f(t)+f(t+1)$，必须同时保证 $t\in[-3,3]$ 和 $t+1\in[-3,3]$；得到普通不等式解集后，还要与这些限制取共同部分。

\medskip
\textbf{\textcolor{CTrap}{考前快速自查}}

\textbf{奇还是偶？严格增还是严格减？定义域是否含 $0$ 且是左右对称的一个区间？代入的自变量都在定义域内吗？}四项都清楚，再使用公式。
\end{trapbox}
